\documentclass[UTF8,11pt]{ctexart}
\usepackage[a4paper,margin=2.35cm]{geometry}
\usepackage{amsmath,amssymb,mathtools}
\usepackage{booktabs,longtable,array}
\usepackage{enumitem}
\usepackage{xcolor}
\usepackage[hidelinks]{hyperref}
\usepackage{url}

\newcommand{\dd}{\,\mathrm d}
\newcommand{\E}{\mathbb E}
\newcommand{\sgn}{\operatorname{sgn}}
\newcommand{\Normal}{\mathcal N}
\newcommand{\cmark}{\textcolor{green!45!black}{\ensuremath{\checkmark}}}
\newcommand{\wmark}{\textcolor{orange!80!black}{\ensuremath{\triangle}}}
\newcommand{\xmark}{\textcolor{red!70!black}{\ensuremath{\times}}}

\title{均值回复与混合 Copula 理论模型：解、推导与适用条件审计}
\author{可复核数学审计}
\date{2026年8月9日}

\begin{document}
\maketitle

\begin{abstract}
本文逐项核验 \path{../meanreversion_article_main.tex} 中的均值回复定义、
Gaussian--Gamma 构造、OU/CIR 单调变换、混合 Copula 构造及债券定价
PDE。核验分为三个层次：代数残差是否为零；随机过程结论所需的可积性、
边界和真鞅条件是否齐备；模型能否产生唯一且有限的观测转移密度。结论是：
仿射 OU、指数 OU、仿射 CIR、二次 CIR 和指数 CIR 的显式系数解在定义域内
代数正确；二次 CIR 在标准参数 $k>0$ 下不能在无限时域保持全局单调；
OU 特殊函数解虽满足 PDE，却因关于状态为偶函数而不能作为全实轴上的单调
变换；CIR 的 Kummer 分离解正确但分支、归一化和可积性条件没有给全；
文中另一个把 CIR 方程化为 $u_\tau=zu_{zz}$ 的变量替换不能通过反代；
Gamma 两个例子只验证了均值点处的零点条件，并未证明全状态空间的符号条件。
混合 Copula 的 Poisson 重置构造和非局部债券 PDE 在补充标准条件后成立。
\end{abstract}

\section{审计范围、符号与判定标准}

为避免原文中同一个 $\theta$ 同时表示基础扩散的长期水平与变换后过程的
目标均值，本文用 $\theta_r$ 表示基础扩散参数，用 $\vartheta$ 表示变换后
过程的目标均值；核对原文公式时再令 $\theta_r=\vartheta=\theta$。标准 CIR
写为
\begin{equation}
 \dd r_t=k(\theta_r-r_t)\dd t+\sigma\sqrt{r_t}\dd W_t,
 \qquad r_t\in[0,\infty).
 \label{eq:cir}
\end{equation}
原文的变换 PDE 在 $\alpha(t)=-d$ 时为
\begin{equation}
 f_t+\mathcal L f+d(f-\vartheta)=0,
 \label{eq:master-pde}
\end{equation}
其中 $\mathcal L$ 是基础扩散的生成元。

本文使用以下判定：
\begin{enumerate}[label=(\roman*)]
 \item \textbf{代数正确}：把候选解反代 ODE/PDE 后恒等残差为零；
 \item \textbf{随机过程有效}：还要求 It\^o 积分为真鞅、所需矩有限且
       状态空间不被破坏；
 \item \textbf{观测密度有效}：对 $Y_t=f(t,r_t)$，还要求每个观测时刻
       $x\mapsto f(t,x)$ 在基础状态空间上单射、观测落在其像集内，且
       Jacobian 有限。此时
 \begin{equation}
 q_Y(y_t\mid y_s)=p_r\!\left(f^{-1}(t,y_t)\mid f^{-1}(s,y_s)\right)
 \left|\frac{\partial f^{-1}(t,y_t)}{\partial y_t}\right|.
 \label{eq:change-density}
 \end{equation}
\end{enumerate}

\section{二次 CIR 中 $k<0$ 的精确结论}

原文二次变换为
\[
 f(t,x)=A(t)x^2+B(t)x+C(t),
\]
其中
\begin{align}
 A(t)&=A_0e^{(2k-d)t},\\
 B(t)&=e^{(k-d)t}\left[B_0+
 \eta\bigl(1-e^{kt}\bigr)\right],
 &\eta&=\frac{A_0(\sigma^2+2k\theta_r)}{k}.
 \label{eq:B-factor}
\end{align}
因为
\[
 f_x(t,x)=2A(t)x+B(t),\qquad x\geq0,
\]
当 $A_0>0$ 时，$x\mapsto f(t,x)$ 在 $x\geq0$ 上单调不减的充要条件是
$B(t)\geq0$。

\subsection{$k>0$：标准 CIR 下只能做有限时域变换}

若 $k>0,\theta_r\geq0,\sigma>0,A_0>0$，则 $\eta>0$，而
$1-e^{kt}\to-\infty$。所以对任意有限 $B_0$，存在有限 $t_\star$ 使
$B(t)<0$。因此原文的 $A(t)>0,B(t)\geq0$ 不能在 $t\in[0,\infty)$ 上
同时成立。固定有限样本窗 $0\leq t\leq T$ 时，单调性等价于
\begin{equation}
 B_0\geq \eta(e^{kT}-1).
 \label{eq:finite-window-monotonicity}
\end{equation}
所以二次变换可以作为有限时域模型，但不能按原文表述直接作为全时域模型。

\subsection{$k<0$：映射可以单调，但基础过程不再是标准 CIR}

当 $k<0$ 时，$e^{kt}$ 从 $1$ 单调下降到 $0$。式 \eqref{eq:B-factor}
中括号的取值位于 $B_0$ 与 $B_0+\eta$ 之间，故全时域单调条件恰为
\begin{equation}
 A_0>0,\qquad \min\{B_0,B_0+\eta\}\geq0.
 \label{eq:knegative-monotonicity}
\end{equation}
所以，若只问“二次映射的单调性假定能否满足”，答案是\textbf{可以}。
但若“模型假定”包括原文的标准均值回复 CIR、Gamma 稳态和正的长期水平，
答案是\textbf{不可以}。

由 \eqref{eq:cir} 得
\begin{equation}
 \E[r_t\mid r_s]
 =\theta_r+(r_s-\theta_r)e^{-k(t-s)}.
 \label{eq:cir-mean}
\end{equation}
标准 CIR 的均值回复、非退化 Gamma 稳态及文中所用参数化要求
$k>0,\theta_r>0,\sigma>0$。若强行令 $k<0$：
\begin{enumerate}
 \item $\theta_r>0$ 时，边界漂移 $k\theta_r<0$，在 $r=0$ 处把过程推出
       非负状态空间，不能作为全局非负 CIR；
 \item $\theta_r=0$ 时可得到超临界分枝型平方根扩散，但
       $\E r_t=r_s e^{|k|(t-s)}$，不均值回复，且转移律一般含边界原子；
 \item $\theta_r<0$ 时 $k\theta_r>0$，可定义非负的\emph{超临界平方根
       扩散}，但漂移在整个 $x\geq0$ 上向上，式 \eqref{eq:cir-mean}
       发散，也不存在原文使用的 Gamma 稳态。若仍共用参数，目标水平还会
       变成负数。
\end{enumerate}
在第三种放宽解释下，时间依赖变换本身仍可能使 $Y_t=f(t,r_t)$ 满足以
$d>0$ 为速率的条件均值回复；但这应重新命名为“基于超临界平方根扩散的
变换模型”，并将基础参数 $\theta_r$ 与目标均值 $\vartheta$ 分开。其理论、
转移密度、稳态叙述和实证参数空间均不同于原文，不能仅把现有 CIR 估计中的
$k$ 改成允许负值。

\section{统一验证引理}

对任意满足 \eqref{eq:master-pde} 的 $C^{1,2}$ 函数，It\^o 公式给出
\begin{equation}
 \dd Y_t=-d(Y_t-\vartheta)\dd t+
 \sigma(r_t)f_x(t,r_t)\dd W_t.
 \label{eq:transformed-sde}
\end{equation}
若局部鞅项在任意有限区间上是真鞅，则
\begin{equation}
 \E[Y_t-\vartheta\mid\mathcal F_s]
 =e^{-d(t-s)}(Y_s-\vartheta).
 \label{eq:conditional-mr}
\end{equation}
因此原文变换方法的核心正确；但 $f\in C^{1,2}$ 本身只保证局部 It\^o
分解，并不足以保证真鞅、有限一阶矩或 $t\to\infty$ 时可交换极限。至少
还要给出增长条件，或直接验证随机积分的一致可积性。

原文一般扩散部分写“$\alpha(t)$ 是正常数 $\alpha>0$”并推出
$\mu(x)=\alpha(x-\theta)$，与前文假定 $\alpha(t)<0$ 冲突。正确写法应为
$\alpha<0$，或令 $d=-\alpha>0$，从而 $\mu(x)=d(\theta-x)$。

\section{基础模型与一般命题}

\subsection{OU 与 CIR}

对 $\dd r_t=k(\theta_r-r_t)\dd t+\sigma\dd W_t$，当 $k>0$ 时
\begin{align*}
 r_t\mid r_s&\sim \Normal\left(
 \theta_r+(r_s-\theta_r)e^{-k\Delta},
 \frac{\sigma^2}{2k}(1-e^{-2k\Delta})\right),\\
 r_\infty&\sim\Normal\left(\theta_r,\frac{\sigma^2}{2k}\right).
\end{align*}
两式均正确。对 CIR，令
\[
 c_\Delta=\frac{\sigma^2(1-e^{-k\Delta})}{4k},\quad
 \nu=\frac{4k\theta_r}{\sigma^2},\quad
 \lambda_\Delta=
 \frac{4ke^{-k\Delta}r_s}{\sigma^2(1-e^{-k\Delta})}.
\]
标准参数下
\[
 r_t\mid r_s\overset d=c_\Delta\chi'^2_\nu(\lambda_\Delta),\qquad
 r_\infty\sim\Gamma\!\left(\frac{2k\theta_r}{\sigma^2},
 \frac{\sigma^2}{2k}\right).
\]
非负弱解需要 $k\theta_r\geq0$；原文的均值回复和非退化稳态还要求
$k>0,\theta_r>0$。Feller 条件 $2k\theta_r\geq\sigma^2$ 保证零边界
不可达，但不是在正观测点写出转移密度的必要条件。

\subsection{定义、生成元与鞅变换}

原文用
\[
 \lim_{h\downarrow0}\frac{\E[X_t\mid X_{t-h}=x]-x}{h}
\]
定义局部漂移。这是从 $t-h$ 到 $t$ 的前向漂移，只是以终点 $t$ 记时。
若要与时变生成元 $A_t$ 完全一致，应统一为
$\lim_{h\downarrow0}\{\E[X_{t+h}\mid X_t=x]-x\}/h$，或者明确使用
原文假设 \texttt{hyp2 relation} 把左右极限接合。当前定义中的
$x\in\operatorname{Ran}(X_t)$ 也应相应改为起点时刻的状态空间。

条件
\[
 \E[X_t-\vartheta\mid X_s]
 =R_{s,t}(X_s-\vartheta),\qquad
 R_{s,t}=\exp\left\{\int_s^t\alpha(u)\dd u\right\},
\]
在 $\alpha<0,R_{s,t}\to0$ 和可积条件下确实推出均值回复。变换
$M_t=(X_t-\vartheta)/R_{0,t}$ 为鞅的代数也正确。需要补充：条件期望应
以 $\mathcal F_s$ 或 Markov 性明确连接；$X_t$ 必须可积；
$R_{0,t}>0$ 使变换可逆，因而保留 Markov 性与 Copula。

\subsection{Gaussian Copula 生成元}

令标准化 OU 满足
$\dd Z_t=-aZ_t\dd t+\sqrt{2a}\dd W_t$，并令 $U_t=\Phi(Z_t)$。It\^o
公式给出
\[
 \mathcal A g(u)=-2a z\phi(z)g'(u)+a\phi(z)^2g''(u),
 \qquad z=\Phi^{-1}(u),
\]
与原文式 (OU generator) 一致。若 $Q=F^{-1}$ 且
$g_F(z)=Q(\Phi(z))$，则
\[
 \mathcal A Q(\Phi(z))=a\{g_F''(z)-zg_F'(z)\}.
\]
所以原文稳态边缘命题给出的全局符号条件确为充分条件。问题在于原文要求
$F^{-1}\in C_0^2([0,1])$，而 Normal 和 Gamma 等无界分布的分位数及其
导数在端点通常发散。应改用开区间上的扩展生成元定义域，并增加端点可积或
局部化论证。

\section{Gaussian Copula--Gamma 两个例子}

设 $Q_\alpha$ 为 $\Gamma(\alpha,\theta_s)$ 的分位数，
$x_\alpha=F_{\alpha,\theta_s}(\alpha\theta_s)$，
$z_\alpha=\Phi^{-1}(x_\alpha)$。原文推导出的
\begin{equation}
 \Gamma(\alpha)\left(\frac e\alpha\right)^\alpha
 \phi(z_\alpha)-2z_\alpha=0
 \label{eq:gamma-root}
\end{equation}
代数上正确：它等价于在边缘均值处令 $\mathcal A Q_\alpha=0$。但是
均值回复要求
\begin{equation}
 \sgn\{\mathcal A Q_\alpha(u)\}
 =\sgn\{\alpha\theta_s-Q_\alpha(u)\}
 \quad\text{对所有 }u\in(0,1),
 \label{eq:global-gamma-sign}
\end{equation}
而式 \eqref{eq:gamma-root} 只检查一个点。原文没有证明
$\mathcal A Q_\alpha$ 只有一个零点，也没有证明零点两侧的符号，所以把
\eqref{eq:gamma-root} 称为充分条件不严格。高精度数值扫描显示式
\eqref{eq:gamma-root} 两边移项后的残差在预先设定的
$\alpha\in[10^{-2},10^3]$、2,000 个对数网格点上保持正值；扩大范围的附加
高精度检查也只显示残差趋近于零。它们都是“未发现有限根”的数值证据，
不是解析不存在证明。

时变尺度例子的 ODE
\[
 \frac{\theta_s'(t)}{\theta_s(t)}
 =\frac{aG(x_t)}{G_\alpha(x_t)}
\]
也只是在移动的目标均值分位点令漂移为零，仍未验证全局符号条件。更强的
必要条件是：若 $\theta_s(t)\to\theta_s^*\in(0,\infty)$，则自治 ODE
右端的极限必须为零，因而端点仍必须满足稳态根条件
\eqref{eq:gamma-root}。所以 Case 2 不能绕开 Case 1 的端点可行性问题。

\section{OU 变换逐项核验}

以下令 $k>0,d>0$，基础 OU 的长期均值为零。

\subsection{仿射 OU}

\[
 f(t,x)=A_0e^{(k-d)t}x+B_0e^{-dt}
\]
反代原文 PDE 后 $x$ 项与常数项残差均为零。$A_0>0$ 保证严格单调；
在基础 OU 平稳时，原文给出的 Normal 均值与方差正确。若 $k>d$，方差
随日历时间增长并不妨碍一阶条件均值回复，但说明 $Y_t$ 不是平稳过程。

\subsection{指数 OU}

原文
\begin{align*}
 A(t)&=A_0e^{kt},\\
 B(t)&=B_0-\frac{\sigma^2A_0^2}{4k}(e^{2kt}-1)-dt,\\
 C(t)&=C_0e^{-dt}
\end{align*}
的三个 ODE 残差均为零。$A_0>0$ 保证单调，Normal 指数矩保证有限时域
可积。若要支持集严格位于正实数，还需 $C_0\geq0$；仅有 $A(t)>0$ 不能
推出这一点。原文的对数正态密度和 Jacobian 正确，支持为 $y>C(t)$。

\subsection{OU 特殊函数}

令 $z=kx^2/\sigma^2$，分离解
\[
 f(t,x)=e^{-ct}J\!\left(\frac{c-d}{2k},\frac12;z\right)
\]
确实把 PDE 化为 Kummer 方程
$zJ_{zz}+(\tfrac12-z)J_z-pJ=0$，所以它是正确的 PDE 解。但是它只依赖
$x^2$，故对每个非零 $x$ 有 $f(t,x)=f(t,-x)$。任何非恒定分支都不可能
在 OU 的全状态空间上单射或全局单调。因此该项不能按原文命题直接宣称为
“OU 的单调递增变换”，也不能用单根 Jacobian 构造观测密度；若引入两个
逆像分支，Copula 保持结论也随之改变。

\section{CIR 变换逐项核验}

\subsection{仿射 CIR}

\begin{align*}
 A(t)&=A_0e^{(k-d)t},\\
 B(t)&=\vartheta\{1-A_0e^{(k-d)t}\}+B_0e^{-dt}
\end{align*}
在原文令 $\vartheta=\theta_r=\theta$ 时两个 ODE 残差为零。
$A_0>0,d>0$ 和标准 CIR 条件足以得到单调、可积及
\eqref{eq:conditional-mr}。观测支持为 $y\geq B(t)$，因此实证中仍须
逐点检查支持。

\subsection{二次 CIR}

原文 $A,B,C$ 三个显式系数反代后全部为零残差；其代数解正确，且
$k\ne0$ 是隐含条件。严格问题是单调性：标准 $k>0$ 下只能满足有限窗条件
\eqref{eq:finite-window-monotonicity}。逆变换在 $A>0,B\geq0$ 时取唯一
非负根
\[
 x=\frac{2(y-C)}{B+\sqrt{B^2+4A(y-C)}},\qquad
 \left|\frac{\partial x}{\partial y}\right|=\frac{1}{2Ax+B},
\]
并要求 $y\geq C(t)$。当 $B(t)=0,y=C(t)$ 时 Jacobian 奇异，连续密度的
端点处理需单独说明。

\subsection{指数 CIR}

Riccati 方程的解
\begin{align*}
 A(t)&=\frac{2k}{\sigma^2+(2k/A_0-\sigma^2)e^{-kt}},\\
 C(t)&=\vartheta+(C_0-\vartheta)e^{-dt}
\end{align*}
以及原文 $B(t)$ 的对数表达式反代后均为零残差。对标准
$k>0,A_0>0$，分母为正且 $A(t)>0$。若基础 CIR 从平稳 Gamma 分布启动，
为使 $\E e^{A_0r_0}<\infty$，还需
\begin{equation}
 0<A_0<\frac{2k}{\sigma^2}.
 \label{eq:cir-exp-integrability}
\end{equation}
若从确定初值启动，则每个有限时点的临界指数更宽，但仍须说明初始分布。
观测支持为 $y>C(t)+e^{B(t)}$。公式正确不代表稳定币价格必然覆盖该支持；
当前实证失败主要属于这一层，而不是 ODE 解错误。

若纯代数上讨论 $k<0$，原文把两个对数分开写会出现负数对数；可将它们合并
为一个正比值的对数后继续。但这对应超临界平方根扩散，不是标准 CIR。

\subsection{CIR 的 Kummer 分离解}

令
\[
 z=\frac{2kx}{\sigma^2},\qquad
 p=\frac{c-d}{k},\qquad q=\frac{2k\theta_r}{\sigma^2}.
\]
则
\[
 f(t,x)=\vartheta+e^{-ct}J(p,q;z)
\]
反代后恰好要求
\[
 zJ_{zz}+(q-z)J_z-pJ=0.
\]
所以 Kummer 分离推导正确。问题是 $J=C_1M+C_2U$ 并未指定分支、常数、
边界条件和归一化，不能据此得到唯一模型。并非所有分支都失败：当
$k>0,q>0,p>0,C_1>0,C_2=0$ 时
\[
 \frac{\partial}{\partial z}M(p,q,z)
 =\frac pq M(p+1,q+1,z)>0,
\]
故存在全局递增的 PDE 分支。但是
$M(p,q,z)\sim e^zz^{p-q}\Gamma(q)/\Gamma(p)$，它对平稳
$z\sim\Gamma(q,1)$ 不可积；从确定初值启动时每个有限时点可积，并可通过
半群特征函数关系获得 \eqref{eq:conditional-mr}。严格结论是：
“该族没有被充分规定，某些分支在某些初始条件下有效”，而不是“特殊函数
模型在数学上不可能”。当前实证中预先给定的有限分支没有同时通过单调、支持
和数值稳定检查，只能说明这些分支失败。

\subsection{声称化为 $u_\tau=zu_{zz}$ 的第二个变量替换}

原文随后给出
\[
 f=\vartheta+e^{-dt}u(\tau,z)
 \exp\!\left\{\sqrt{2k/\sigma^2}\,x-k\theta_r t\right\}
\]
及相应 $z,\tau$，并声称得到 $u_\tau=zu_{zz}$。该变换不正确。
例如取其声称的解 $u=z$，反代原 CIR PDE；除去非零公共指数因子后，残差为
\begin{align*}
 R(x)={}&k\theta_r+
 \left[2\sqrt{2k\sigma^2}-k-k\theta_r
 +k\theta_r\sqrt{2k/\sigma^2}\right]x\\
 &+k\left(1-\sqrt{2k/\sigma^2}\right)x^2.
\end{align*}
一般不为零；例如 $k=2,\sigma=1,\theta_r=1,x=1$ 时 $R=4$。因此其后
列出的三组 $u$ 虽然各自满足 $u_\tau=zu_{zz}$，却没有被这个错误变换证明
为原 CIR PDE 的解。

一个直接的正确降维是令 $v=f-\vartheta=e^{-dt}u(\tau,z)$、
\[
 z=xe^{kt},\qquad
 \tau=\frac{\sigma^2}{2k}(1-e^{kt}).
\]
反代得到
\[
 u_\tau=zu_{zz}+\frac{2k\theta_r}{\sigma^2}u_z,
\]
仍保留一阶项；只有特殊参数退化时才化为原文所写的无一阶项方程。

\section{混合 Copula、Poisson 构造与债券 PDE}

设基础 Markov 转移核为 $P_{s,t}$，平稳边缘核为
$\Pi(y,\dd z)=F(\dd z)$，并定义
\[
 P^{\mathrm{mix}}_{s,t}=a_{s,t}P_{s,t}+(1-a_{s,t})\Pi,\qquad
 a_{s,t}=e^{-\int_s^tc(u)\dd u}.
\]
因为 $a_{s,u}=a_{s,t}a_{t,u}$、$P_{s,t}P_{t,u}=P_{s,u}$，且
$P\Pi=\Pi P=\Pi$，直接相乘得到
\[
 P^{\mathrm{mix}}_{s,t}P^{\mathrm{mix}}_{t,u}
 =P^{\mathrm{mix}}_{s,u}.
\]
因此混合 Copula 的一致性结论正确，前提是原文的“consistent”指
Markov Chapman--Kolmogorov 一致性，并且基础核保持同一个平稳边缘。

对独立平稳副本 $Y^{(i)}$ 与强度 $\lambda$ 的 Poisson 过程，
$Y_t=Y_t^{(N_t)}$ 在区间 $(s,t]$ 无跳时沿原核演化，至少一次跳时终值来自
独立平稳边缘。因此
\[
 P^{\mathrm{mix}}_\Delta(y,\dd z)
 =e^{-\lambda\Delta}P_\Delta(y,\dd z)
 +(1-e^{-\lambda\Delta})F(\dd z),
\]
原文的概率构造和 Copula 权重正确。若基础过程的平稳均值为 $\vartheta$，
混合过程的局部漂移为
\[
 \mu_{\mathrm{mix}}(y)=\mu_0(y)+\lambda(\vartheta-y),
\]
故保留均值回复符号。

以 OU 为基础时，混合过程生成元为
\[
 \mathcal Gg(y)=k(\vartheta-y)g_y+\frac12\sigma^2g_{yy}
 +\lambda\left\{\int g(z)F(\dd z)-g(y)\right\}.
\]
Feynman--Kac 方程 $g_t+\mathcal Gg-yg=0,\ g(T,y)=1$ 与原文债券 PDE
一致。严格版本需说明风险中性测度、指数泛函可积性，以及采用经典解还是
黏性解；这些是定理条件的缺失，不是 PDE 符号错误。

\section{汇总结论与修正级别}

\begin{longtable}{p{0.25\textwidth}p{0.16\textwidth}p{0.50\textwidth}}
\toprule
对象 & 判定 & 核验结论\\
\midrule
\endhead
基础 OU/CIR & \cmark 条件成立 & 解、条件均值和标准转移律正确；需显式写出参数与边界条件。\\
Gaussian Copula 生成元 & \cmark & It\^o 推导正确；分位数函数的生成元定义域需重写。\\
Gamma Case 1 & \wmark 不充分 & 标量方程只是一点处必要条件，没有证明全局符号。\\
Gamma Case 2 & \wmark 不充分 & ODE 只是一点处必要条件，且收敛端点必须满足 Case 1。\\
鞅与一般扩散变换 & \wmark 条件成立 & 核心 PDE 正确；“当且仅当”需要真鞅、增长和可积条件。\\
仿射 OU & \cmark & 显式解正确，补 $A_0>0,d>0$。\\
指数 OU & \cmark 条件成立 & 显式解和密度正确；正支持需 $C_0\geq0$。\\
OU 特殊函数 & \xmark 非单射 & PDE 解正确，但偶函数结构导致全实轴非单射。\\
仿射 CIR & \cmark & 显式解正确，需标准 CIR、单调和支持条件。\\
二次 CIR & \wmark 仅有限时域 & 系数解正确；标准 $k>0$ 下全时域单调不可能。\\
指数 CIR & \cmark 条件成立 & Riccati 解正确；必须补指数矩和随时间支持条件。\\
CIR Kummer 分离解 & \wmark 未充分规定 & 分离推导正确，存在单调分支；需指定分支、归一化、初始分布和可积性。\\
CIR 第二变量替换 & \xmark 错误 & 反代残差非零，其后三个 $u$ 不能据此视为 CIR 解。\\
混合 Copula/Poisson 重置 & \cmark & 核组合和概率构造正确，依赖平稳与独立假设。\\
混合 OU 债券 PDE & \cmark 条件成立 & 生成元与 Feynman--Kac 符号正确，需补解存在性条件。\\
\bottomrule
\end{longtable}

对实证的直接影响是：
\begin{enumerate}
 \item 二次 CIR 的现有失败诊断应保留；它反映标准 $k>0$ 下的全时域单调
       障碍。允许 $k<0$ 会改变基础模型类别，不能作为同一模型的简单扩展。
 \item 指数 CIR 的理论 ODE 解可以保留，其实证失败应表述为支持、数值条件
       与识别失败，而不是解析解错误。
 \item CIR 特殊函数模型应从“结构上不可能”改为“当前预设分支未通过”；
       若继续估计，应先固定有限分支族及归一化，再证明相应初始分布下可积。
 \item OU 特殊函数分支是真正的全局单射失败；若改成双分支隐藏状态模型，
       已不再是原文的单调 Copula 保持变换。
 \item Gamma 两个模型不应进入“已证明均值回复”的集合，只能作为必要条件
       诊断，直到全局符号与端点存在性被证明。
 \item 混合 Copula 的核心 Poisson 重置模型及债券 PDE 没有发现代数错误；
       其理论主线可以保留并补全条件。
\end{enumerate}

\section{可复核材料与文档完整性}

符号核验脚本为 \path{scripts/15_verify_theory_derivations.py}；机器可读结论为
\path{output/tables/theory_derivation_audit.csv}；零残差和反例为
\path{output/models/theory_symbolic_checks.json}。

原理论文第 691 行开始的 “Proofs” 章节为空，且正文主动引用但未在当前文件
定义的关键标签包括 \texttt{copula-generator}、\texttt{time change copula}、
\texttt{time change generator}、\texttt{dynamic martingale condition}、
\texttt{E1} 和 \texttt{diffusion copula}。因此即使本审计补做了主要推导，
当前 \path{../meanreversion_article_main.tex} 仍不是一份自包含、可逐定理
核查的完整理论稿。正式投稿前应补入这些定义、定理及全部证明，并按本报告
修正条件。

\end{document}
